\chapter{Rozmaitości Seiferta i grafy gwiaździste}
\label{ch:tg-seifert}

Z grafów plumbingowych odczytujemy formy przecięcia.
Przy odpowiednich założeniach dotyczących grafu możemy
ponadto rozpoznać rozwarstwienie jego brzegu na okręgi.
To prowadzi do \emph{rozmaitości Seiferta}. Zaczniemy od lokalnego
modelu, a następnie wyprowadzimy prezentację grupy fundamentalnej,
formułę na pierwszą homologię i związek z ułamkami łańcuchowymi.
W tym rozdziale rozpatrujemy zamknięte, spójne,
zorientowane trójrozmaitości Seiferta o orientowalnej
orbipowierzchni bazowej. Torusy pełne i rozmaitości
z brzegiem pojawiają się jako części pomocnicze
konstrukcji. Nie omawiamy tutaj odmian z nieorientowalną
bazą ani z brzegiem lustrzanym.

\section{Włókna zwykłe i wyjątkowe: model lokalny}
\label{sec:tg-seifert-model}

\begin{definicja}[Rozwarstwienie Seiferta]
Niech $M$ będzie zamkniętą, spójną, zorientowaną
gładką trójrozmaitością.
\emph{Rozwarstwieniem Seiferta} nazywamy rozkład $M$ na rozłączne
okręgi, dla którego każdy okrąg ma nasycone otoczenie dyfeomorficzne
z torusem pełnym $V_{\alpha,q}=(D^2\times S^1)/\mathbb Z_\alpha$,
z rozkładem na obrazy zbiorów $\{z\}\times S^1$.
Przy $\alpha\geq1$ grupa cykliczna działa według wzoru
\[
(z,u)\longmapsto(\zeta z,\zeta^q u),\qquad
\zeta=e^{2\pi i/\alpha},\quad \gcd(\alpha,q)=1.
\]
Dla $\alpha=1$ przyjmujemy działanie trywialne. Obraz
$\{0\}\times S^1$ jest \emph{włóknem wyjątkowym krotności $\alpha$},
gdy $\alpha>1$; włókna krotności jeden nazywamy \emph{zwykłymi}.
\end{definicja}

\begin{uwaga}[Dlaczego iloraz jest rozmaitością]
Jeśli potęga generatora o wykładniku $j$ ustala $(z,u)$, to
$\zeta^{qj}u=u$, czyli $\alpha\mid qj$. Ponieważ
$\gcd(\alpha,q)=1$, wynika stąd $\alpha\mid j$:
element ustalający punkt jest tożsamością.
Działanie jest więc wolne. Dla $\alpha>1$ wybieramy
$k\in\mathbb Z$ takie, że $qk\equiv1\pmod\alpha$.
Odwzorowanie
\[
F(z,u)=(zu^{-k},u^\alpha)
\]
jest niezmiennicze względem działania generatora:
$\zeta z(\zeta^qu)^{-k}=\zeta^{1-qk}zu^{-k}=zu^{-k}$,
a $(\zeta^qu)^\alpha=u^\alpha$.
Indukuje więc gładkie odwzorowanie ilorazu na
$D^2\times S^1$. Dla $(v,w)\in D^2\times S^1$
wybieramy pierwiastek $u^\alpha=w$ i kładziemy
$z=vu^k$. Zmiana pierwiastka na $\zeta^j u$
zmienia parę na $(\zeta^{jk}z,\zeta^ju)$,
czyli na obraz przez $(jk)$-tą potęgę generatora,
gdyż $qk\equiv1\pmod\alpha$.
Odwrotność jest zatem dobrze określona,
a lokalne gładkie wybory pierwiastka pokazują
jej gładkość. Dla $\alpha=1$ iloraz jest
bezpośrednio produktem. Otrzymujemy torus pełny,
bez osobliwości rozmaitości. Przestrzeń włókien
to iloraz $D^2/\mathbb Z_\alpha$ przez obroty,
z lokalną współrzędną topologiczną $z^\alpha$.
Dla $\alpha>1$ jedynie środek ma nietrywialną
grupę izotropii; odpowiada mu kąt stożkowy
$2\pi/\alpha$ w standardowym modelu metrycznym
ilorazu dysku przez obroty. Parametr $q$ wpływa na sposób sklejenia
włókna zwykłego z meridianem torusa pełnego.
\end{uwaga}

Przestrzeń włókien jest dwuwymiarową \emph{orbipowierzchnią}
$B$. W rozpatrywanym przypadku ma orientowalną powierzchnię
podstawową rodzaju $g$ i skończenie wiele punktów stożkowych
rzędów $\alpha_1,\ldots,\alpha_r$. Oznaczamy ją
$B(g;\alpha_1,\ldots,\alpha_r)$. Skończoność
zbioru punktów stożkowych wynika ze zwartości:
każde lokalne otoczenie w bazie zawiera co najwyżej
jeden taki punkt, a skończone podpokrycie tych
otoczeń obejmuje całą bazę. Topologiczna
przestrzeń podkładowa orbipowierzchni nie pamięta rzędów
punktów stożkowych; dlatego samej powierzchni $B$ nie wolno
utożsamiać z pełnymi danymi rozwarstwienia.

\begin{center}
\begin{tikzpicture}[scale=.91,>=Stealth,line cap=round]
 \fill[tgBlue!9] (-3.6,-.1) ellipse (1.65 and .72);
 \draw[tgBlue,line width=1.4pt] (-3.6,-.1) ellipse (1.65 and .72);
 \draw[tgNavy!50,dashed] (-3.6,-.1) ellipse (.66 and .72);
 \foreach \x in {-4.1,-3.1}{\draw[tgTeal,line width=1.1pt] (\x,-.65) arc[start angle=-90,end angle=90,x radius=.19,y radius=.53];}
 \draw[tgCoral,line width=2.6pt] (-3.6,-.08) ellipse (.46 and .30);
 \node[font=\scriptsize,text=tgCoral] at (-3.6,-1.10) {włókno centralne};
 \draw[tgTeal,-{Stealth[length=2.8mm]},line width=1.6pt] (-1.3,-.1)--(1.05,-.1);
 \node[font=\small,text=tgTeal] at (-.1,.25) {iloraz};
 \fill[tgBlue!8] (3,-.1) ellipse (1.2 and .75);
 \draw[tgBlue,line width=1.4pt] (3,-.1) ellipse (1.2 and .75);
 \fill[tgCoral] (3,-.1) circle (.095);
 \node[font=\scriptsize,text=tgNavy] at (3,-1.11) {punkt stożkowy rzędu $\alpha$};
 \node[font=\scriptsize,text=tgNavy] at (0,-1.66)
 {Model lokalny: okręgi w torusie pełnym i odpowiadający im dysk bazowy.};
\end{tikzpicture}
\end{center}

\section{Niezmienniki Seiferta jako dane sklejenia}
\label{sec:tg-seifert-gluing}

Aby uniknąć niejednoznaczności znaków, zdefiniujemy
niezmienniki Seiferta przez konkretne nachylenia chirurgii.
Niech $B_{g,r+1}$ będzie zwartą, spójną, zorientowaną
powierzchnią rodzaju $g$ z $r+1$ brzegami, oznaczonymi
$\gamma_0,\gamma_1,\ldots,\gamma_r$.
Dla krzywych na powierzchni wybieramy punkt bazowy i drogi
do brzegów tak, aby w $\pi_1(B_{g,r+1})$ zachodziła relacja
\[
\gamma_0\gamma_1\cdots\gamma_r
\prod_{j=1}^g[a_j,b_j]=1.
\]
Na każdym torusie brzegowym $B_{g,r+1}\times S^1$ mamy
klasę krzywej $\gamma_i$ oraz klasę włókna $h$.

\begin{definicja}[Zorientowane dane Seiferta]
Niech $g\geq0$, $b\in\mathbb Z$, $r\geq0$,
$\alpha_i\geq2$ i $0<\beta_i<\alpha_i$ oraz
$\gcd(\alpha_i,\beta_i)=1$ dla $1\leq i\leq r$.
Przez
$M(g;b;(\alpha_1,\beta_1),\ldots,(\alpha_r,\beta_r))$
oznaczamy orientowaną rozmaitość powstającą z
$B_{g,r+1}\times S^1$ przez dołączenie torusów pełnych
w taki sposób, aby klasy meridianów dołączanych torusów
były odpowiednio
$[\gamma_0]+b[h]$ oraz
$\alpha_i[\gamma_i]+\beta_i[h]$ dla $1\leq i\leq r$.
W tej konwencji \emph{wymierną liczbą Eulera} nazywamy
$e_{\mathbb Q}(M):=-(b+\sum_{i=1}^r\beta_i/\alpha_i)$.
\end{definicja}

Wypełnienie torusa odpowiadającego $\gamma_0$ jest zwykłe
(tj. nie tworzy włókna wyjątkowego) i koduje liczbę
całkowitą $b$. Pozostałe tworzą włókna wyjątkowe
krotności $\alpha_i$. Aby uzasadnić istnienie
sklejeń i odczyt krotności, na danym torusie
zapiszmy nachylenie meridianu jako
$m=\alpha[\gamma]+\beta[h]$, gdzie dla zwykłego
wypełnienia $\alpha=1$ i $\beta=b$.
Względna pierwszość $\alpha,\beta$ pozwala
wybrać liczby całkowite $p,s$ z
$\alpha s-\beta p=1$. Klasa
$\ell=p[\gamma]+s[h]$ wraz z $m$ tworzy
całkowitą bazę homologii torusa. W tej bazie $[h]=-p m+\alpha\ell$.
Przy orientacjach brzegowych indukowanych przez
orientacje produktowe bazy $(\gamma,h)$ oraz
meridian--długość torusa pełnego są dodatnie.
Sklejenie zorientowanych części musi odwracać
orientację brzegu. Dlatego dodatni meridian
torusa dołączanego wysyłamy na $-m$, a dodatnią
długość na $\ell$. Macierz sklejenia ma wówczas
wyznacznik $-1$; zmiana znaku meridianu nie
zmienia nachylenia chirurgii. Macierz tę realizuje
liniowy dyfeomorfizm torusa. W bazie torusa
dołączanego włókno $h$ ma więc współrzędne
$(p,\alpha)$, a jego stopień względem rdzenia
wynosi $\alpha$.

Lokalny model z poprzedniej sekcji realizuje
to nachylenie: w jego współrzędnych produktowych
$F(z,u)$ krzywa włókna na brzegu ma przy pełnym
obiegu parametru $u$ współrzędne o stopniach
$(-k,\alpha)$. Wybieramy $k=-p$; ponieważ
$\alpha s-\beta p=1$, mamy $\gcd(k,\alpha)=1$
dla $\alpha>1$ i możemy wybrać $q$ odwrotny
do $k$ modulo $\alpha$.
Dla $\alpha=1$ nachylenie $(p,1)$ realizuje
zwykłe rozwarstwienie produktu po zmianie
współrzędnych $(z,u)\mapsto(zu^p,u)$.
Na brzegu oba rozwarstwienia są rodzinami
przesunięć równoległych okręgów o tym samym
pierwotnym wektorze kierunkowym. Liniowe
sklejenie przenosi więc włókna na włókna,
a nie tylko ich klasy homologii. Rozwarstwienie
przedłuża się na wnętrze torusa z krotnością
$\alpha$ włókna rdzeniowego. Są to \emph{dane rozwarstwienia},
a nie tylko niezmienniki samej rozmaitości: niektóre
trójrozmaitości dopuszczają więcej niż jedno rozwarstwienie
Seiferta. Zmiana konwencji orientacji i krzywych brzegowych
może prowadzić do innych zapisów tych samych danych.

\begin{center}
\begin{tikzpicture}[scale=.92,>=Stealth,line cap=round]
 \fill[tgBlue!9] (0,0) ellipse (2.1 and 1.15);
 \draw[tgBlue,line width=1.4pt] (0,0) ellipse (2.1 and 1.15);
 \foreach \x/\col/\lab in {-1.1/tgTeal/$\alpha_1$,0/tgCoral/$\alpha_2$,1.1/tgPurple/$\alpha_3$}{
  \draw[\col,line width=1.2pt] (\x,.1) circle (.33);
  \fill[\col!20] (\x,.1) circle (.29);
  \node[text=\col,font=\small] at (\x,.1) {\lab};
 }
 \node[font=\small,text=tgNavy] at (0,-1.54) {baza $S^2$ z trzema punktami stożkowymi};
 \draw[tgGold,line width=1.4pt,->] (2.45,.35)--(3.4,.35);
 \node[font=\footnotesize,text=tgGold,anchor=west] at (3.46,.35)
 {$S^1$ nad punktami zwykłymi};
\end{tikzpicture}
\end{center}

\section{Grupa fundamentalna: dowód przez van Kampena}
\label{sec:tg-seifert-pi1}

\begin{twierdzenie}[Prezentacja grupy rozmaitości Seiferta]
\label{thm:tg-seifert-presentation}
Niech $M=M(g;b;(\alpha_i,\beta_i)_{i=1}^r)$ będzie
rozmaitością z poprzedniej definicji. Wówczas jej grupa
fundamentalna ma prezentację
\[
\pi_1(M)\cong\Big\langle a_j,b_j,c_i,h\ \Big|\
\begin{array}{l}
[h,a_j]=[h,b_j]=[h,c_i]=1,\\
c_i^{\alpha_i}h^{\beta_i}=1\quad (1\leq i\leq r),\\
c_1\cdots c_r\prod_{j=1}^g[a_j,b_j]=h^b
\end{array}\Big\rangle.
\]
Indeksy $j$ przebiegają $1,\ldots,g$, a $i$ przebiegają
$1,\ldots,r$; dla pustych zbiorów indeksów pomijamy
dane generatory i relacje. Element $h$ jest reprezentowany
przez regularne włókno i należy do centrum grupy.
\end{twierdzenie}

\begin{proof}
Przed wypełnieniem torusów rozważamy produkt
$M_0=B_{g,r+1}\times S^1$. Ponieważ jest produktem,
$\pi_1(M_0)=\pi_1(B_{g,r+1})\times\langle h\rangle$,
gdzie $h$ komutuje z wszystkimi generatorami
powierzchni. Grupa $\pi_1(B_{g,r+1})$ ma
prezentację na generatorach $a_j,b_j,c_0,\ldots,c_r$
z relacją $c_0c_1\cdots c_r\prod_j[a_j,b_j]=1$;
wynika to np. z deformacyjnej retrakcji powierzchni
z brzegiem na odpowiedni graf i wycięcia dysków.

Dołączenie torusa pełnego wzdłuż wybranego
\emph{nachylenia} na torusie brzegowym narzuca relację,
że klasa tego nachylenia jest trywialna w grupie
fundamentalnej wypełnionej rozmaitości. Rzeczywiście
klasa ta jest obrazem meridianu torusa dołączanego,
a meridian ogranicza w nim dysk. Aby zastosować twierdzenie van Kampena, pogrubiamy
obie części w kołnierzu sklejania; ich przecięcie
retraktuje się na spójny torus brzegowy.
Niech $m,\ell$ będą bazą meridian--długość
dołączanego torusa z poprzedniego rachunku.
W grupie fundamentalnej torusa pełnego
$m$ przechodzi na jedynkę, a $\ell$ na generator $t$
grupy nieskończonej cyklicznej.
Jeżeli $i$ oznacza homomorfizm indukowany
włożeniem torusa brzegowego do dotychczasowej części
(o ustalonej drodze do punktu bazowego),
van Kampen daje iloczyn wolny dotychczasowej grupy
z $\langle t\rangle$ z dodatkowymi relacjami
$i(m)=1$ oraz $i(\ell)=t$.
Wystarczą te dwie relacje, ponieważ $m,\ell$
generują grupę fundamentalną torusa, a ich obrazy
już komutują. Eliminując $t$ przez drugą relację,
otrzymujemy dokładnie iloraz dotychczasowej grupy
przez normalne domknięcie $i(m)$.
Nie chodzi zatem o zabicie samej klasy $\gamma_i$,
lecz o zabicie konkretnych klas
$\gamma_0h^b$ oraz $\gamma_i^{\alpha_i}h^{\beta_i}$.
Otrzymujemy dokładnie relacje $c_0h^b=1$ oraz
$c_i^{\alpha_i}h^{\beta_i}=1$.

Z pierwszej relacji mamy $c_0=h^{-b}$. Po podstawieniu
do relacji powierzchniowej uzyskujemy
$c_1\cdots c_r\prod_j[a_j,b_j]=h^b$.
Usunięcie generatora $c_0$ przez przekształcenie
Tietzego daje dokładnie wskazaną prezentację.
Centralność $h$ została odziedziczona z iloczynu
$B_{g,r+1}\times S^1$.
\end{proof}

\begin{wniosek}[Grupa orbipowierzchni jako iloraz]
Dla rozmaitości $M$ z twierdzenia
\ref{thm:tg-seifert-presentation} iloraz
$\pi_1(M)/\langle h\rangle$ ma prezentację
\[
\big\langle a_j,b_j,c_i\mid c_i^{\alpha_i}=1,
\ c_1\cdots c_r\prod_j[a_j,b_j]=1\big\rangle.
\]
Jest to grupa fundamentalna orbipowierzchni
$B(g;\alpha_1,\ldots,\alpha_r)$.
\end{wniosek}
\begin{proof}
Ponieważ $h$ jest centralny, podgrupa cykliczna
$\langle h\rangle$ jest normalna. W prezentacji
ilorazu dodajemy relację $h=1$, która usuwa $h$
ze wszystkich istniejących relacji. Pozostają dokładnie
zapisane generatory i relacje. Aby rozpoznać
grupę orbipowierzchni, usuwamy z jej powierzchni
podkładowej małe dyski wokół punktów stożkowych.
Grupa fundamentalna pozostałej powierzchni ma
prezentację z relacją
$c_1\cdots c_r\prod_j[a_j,b_j]=1$.
Dołączenie dysku stożkowego rzędu $\alpha_i$
w grupie orbipowierzchni narzuca relację
$c_i^{\alpha_i}=1$: w orbipokryciu lokalnym
$D^2\to D^2/\mathbb Z_{\alpha_i}$ dopiero
$\alpha_i$ obiegów brzegu podnosi się do
zamkniętej pętli ograniczającej dysk.
Równoważnie grupa fundamentalna orbipowierzchni
jest z definicji ilorazem grupy części regularnej
przez normalne domknięcie tych potęg meridianów
(por. P.~Scott, \emph{The geometries of 3-manifolds},
sekcja 2). To daje dokładnie tę samą prezentację,
także dla $r=0$.
\end{proof}

\begin{uwaga}[Uwaga na rząd włókna]
Nie zawsze $\langle h\rangle\cong\mathbb Z$:
w rozmaitościach Seiferta o skończonej grupie fundamentalnej
włókno zwykłe może mieć skończony rząd. Dlatego poprawny
zapis sekwencji opisującej grupę jest
$1\to\langle h\rangle\to\pi_1(M)\to
\pi_1^{\mathrm{orb}}(B)\to1$,
a nie bezwarunkowo $1\to\mathbb Z\to\pi_1(M)\to\cdots$.
\end{uwaga}

\section{Pierwsza homologia i wymierna liczba Eulera}
\label{sec:tg-seifert-homology}

\begin{twierdzenie}[Homologia Seiferta nad sferą]
\label{thm:tg-seifert-h1}
Niech $M=M(0;b;(\alpha_i,\beta_i)_{i=1}^r)$,
gdzie $b\in\mathbb Z$, $\alpha_i\ge2$,
$0<\beta_i<\alpha_i$ oraz $\gcd(\alpha_i,\beta_i)=1$.
Grupa $H_1(M;\mathbb Z)$ ma prezentację abelową
z generatorami $c_1,\ldots,c_r,h$ i relacjami
$\alpha_ic_i+\beta_ih=0$ oraz
$c_1+\cdots+c_r-bh=0$.
Jeśli $e_{\mathbb Q}(M)\ne0$, grupa ta jest skończona
rzędu
\[
|H_1(M;\mathbb Z)|=
\left|\alpha_1\cdots\alpha_r
\left(b+\sum_{i=1}^r\frac{\beta_i}{\alpha_i}\right)\right|.
\]
Jeśli $e_{\mathbb Q}(M)=0$, to $H_1(M;\mathbb Z)$ ma
rangę jeden; jej ewentualną część torsyjną wyznacza
postać Smitha macierzy relacji.
\end{twierdzenie}

\begin{proof}
Z twierdzenia Hurewicza w stopniu pierwszym,
ponieważ $M$ jest spójna i lokalnie łukowo spójna,
$H_1(M;\mathbb Z)$ jest abelianizacją $\pi_1(M)$.
Abelianizacja prezentacji z twierdzenia
\ref{thm:tg-seifert-presentation} zabija wszystkie
komutatory. Dla $r=0$ pozostaje jeden generator
$h$ z relacją $bh=0$. Otrzymujemy
$\mathbb Z/b\mathbb Z$ dla $b\ne0$ oraz
$\mathbb Z$ dla $b=0$, zgodnie z tezą przy
konwencji pustego iloczynu równego jeden.
Dalej zakładamy $r\ge1$. Otrzymujemy więc kwadratową macierz
relacji $A$ wymiaru $r+1$: elementy diagonalne
pierwszego bloku to $\alpha_i$, ostatnia kolumna
ma wartości $\beta_i$, a ostatni wiersz to
$(1,\ldots,1,-b)$.
Nad $\mathbb Q$ eliminujemy pierwsze $r$ generatorów.
W blokowej postaci
\[
A=\begin{pmatrix}
\operatorname{diag}(\alpha_1,\ldots,\alpha_r)&\beta\\
\mathbf 1^T&-b
\end{pmatrix}
\]
oznaczamy pierwszy blok diagonalny przez $D$.
Każda $\alpha_i$ jest dodatnia, więc $D$ jest
odwracalny nad $\mathbb Q$. Mnożąc $A$ z lewej
przez macierz
$\begin{psmallmatrix}I&0\\-\mathbf 1^TD^{-1}&1\end{psmallmatrix}$
o wyznaczniku jeden, otrzymujemy macierz
$\begin{psmallmatrix}D&\beta\\0&-b-\mathbf 1^TD^{-1}\beta\end{psmallmatrix}$.
Zatem
\[
\det A=\alpha_1\cdots\alpha_r
\left(-b-\sum_i\frac{\beta_i}{\alpha_i}\right).
\]
Nie jest to na ogół całkowita zmiana bazy,
dlatego nie odczytujemy z niej struktury torsji.
Gdy ta liczba jest niezerowa, homologia jest skończonym
ilorazem $\mathbb Z^{r+1}/A^T\mathbb Z^{r+1}$
(o tej samej wartości wyznacznika), a jej rząd
wynosi $|\det A|$. Jeśli wyznacznik jest zerowy,
pierwszy blok ma nadal pełny rząd $r$, zatem
$\operatorname{rank}_{\mathbb Q}A=r$ i grupa
ilorazowa ma rangę jeden. To dowodzi wszystkich tez.
\end{proof}

\begin{przyklad}[Dane Seiferta nad $S^2(2,3,5)$]
Rozważmy dane $b=-2$ oraz
$(\alpha_1,\beta_1)=(2,1)$,
$(\alpha_2,\beta_2)=(3,2)$,
$(\alpha_3,\beta_3)=(5,4)$.
Wtedy
\[
b+\frac{\beta_1}{\alpha_1}
+\frac{\beta_2}{\alpha_2}
+\frac{\beta_3}{\alpha_3}
=-2+\frac12+\frac23+\frac45=-\frac1{30}.
\]
Z twierdzenia otrzymujemy $|H_1(M;\mathbb Z)|=30/30=1$,
czyli $H_1(M;\mathbb Z)=0$. Nie oznacza to
jeszcze, że $M\cong S^3$: grupa fundamentalna może
być nieprzemienna i doskonała. Pokażemy to poniżej.
\end{przyklad}

\begin{przyklad}[Zerowa liczba Eulera nie wyklucza torsji]
\label{ex:tg-seifert-zero-e-torsion}
Rozważmy rozmaitość Seiferta
\[
M=M(0;-1;(3,1),(3,1),(3,1)).
\]
Jej wymierna liczba Eulera wynosi
$e_{\mathbb Q}(M)=-(-1+3/3)=0$.
Twierdzenie~\ref{thm:tg-seifert-h1}
gwarantuje zatem, że $H_1(M;\mathbb Z)$
ma rangę jeden, nie określa jednak samo
przez się części torsyjnej.

W abelianizacji prezentacji grupy mamy
generatory $c_1,c_2,c_3,h$ i relacje
\[
3c_1+h=3c_2+h=3c_3+h=0,
\qquad c_1+c_2+c_3+h=0.
\]
Z pierwszej relacji otrzymujemy $h=-3c_1$.
Po podstawieniu oraz położeniu
$d_2=c_2-c_1$, $d_3=c_3-c_1$
dostajemy
\[
3d_2=3d_3=0,\qquad d_2+d_3=0.
\]
Eliminując $d_3=-d_2$, widzimy, że
$c_1$ pozostaje wolny, a $d_2$
jest generatorem grupy cyklicznej rzędu
trzy. Zatem
\[
H_1(M;\mathbb Z)\cong
\mathbb Z\oplus\mathbb Z/3\mathbb Z.
\]
Przykład pokazuje konieczność oddzielnego
rachunku całkowitego, np. przez postać normalną
Smitha, gdy $e_{\mathbb Q}=0$.
\end{przyklad}

\begin{wniosek}[Ranga pierwszej homologii dla bazy rodzaju $g$]
\label{cor:tg-seifert-general-h1}
Niech $M=M(g;b;(\alpha_i,\beta_i)_{i=1}^r)$, przy czym
$g\geq0$, $r\geq0$ oraz pozostałe parametry są jak
w definicji danych Seiferta. Wówczas
\[
H_1(M;\mathbb Z)\cong\mathbb Z^{2g}\oplus A,
\]
gdzie $A$ jest grupą abelową o generatorach
$c_1,\ldots,c_r,h$ i relacjach
$\alpha_i c_i+\beta_i h=0$ oraz
$\sum_i c_i-bh=0$.
W szczególności, jeśli $b+\sum_i\beta_i/\alpha_i\ne0$,
to $\operatorname{rank}H_1(M;\mathbb Z)=2g$,
a gdy suma ta jest zerowa, ranga wynosi $2g+1$.
\end{wniosek}
\begin{proof}
Abelianizujemy prezentację z twierdzenia
\ref{thm:tg-seifert-presentation}. Wszystkie komutatory
$[a_j,b_j]$ stają się trywialne. Generatory
$a_1,b_1,\ldots,a_g,b_g$ nie występują w żadnej
pozostałej relacji, wyznaczają więc wolny składnik
$\mathbb Z^{2g}$. Pozostała grupa ma tę samą macierz
relacji, którą obliczyliśmy dla $g=0$.
Jej ranga nad $\mathbb Q$ wynosi zero albo jeden,
zależnie od tego, czy dopełnienie Schura
$-b-\sum_i\beta_i/\alpha_i$ jest niezerowe,
czy równe zeru.
\end{proof}

\begin{przyklad}[Wiązka trywialna nad torusem]
Dla $g=1$, $r=0$ oraz $b=0$ otrzymujemy
$M(1;0;)=T^2\times S^1=T^3$. Z prezentacji grupy
mamy $\pi_1(M)=\langle a,b,h\mid [a,b]=[a,h]=[b,h]=1\rangle
\cong\mathbb Z^3$, w zgodzie z równością
$H_1(T^3;\mathbb Z)\cong\mathbb Z^3$.
\end{przyklad}

\section{Od ramion plumbingowych do rozwarstwienia Seiferta}
\label{sec:tg-seifert-stars}

Niech $\Gamma$ będzie drzewem gwiaździstym z centralną
wagą $b$ oraz $r$ ramionami. W $i$-tym ramieniu wagi
wynoszą $-a_{i1},\ldots,-a_{ik_i}$, licząc od centrum,
gdzie $a_{ij}\ge2$. Załóżmy, że
\[
\frac{\alpha_i}{\beta_i}=
[a_{i1},\ldots,a_{ik_i}]^-
:=a_{i1}-\cfrac1{a_{i2}-\cfrac1{\ddots-\cfrac1{a_{ik_i}}}},
\quad \alpha_i>\beta_i>0,\quad\gcd(\alpha_i,\beta_i)=1.
\]
Dane $(\alpha_i,\beta_i)$ wyznaczają zatem ułamki
pochodzące z poszczególnych ramion.

\begin{stwierdzenie}[Wyznacznik grafu gwiaździstego]
\label{prop:tg-star-determinant}
Przy powyższych założeniach macierz przecięcia
$Q_\Gamma$ spełnia
\[
\big|\det Q_\Gamma\big|=
\left|\prod_{i=1}^r\alpha_i
\left(b+\sum_{i=1}^r\frac{\beta_i}{\alpha_i}\right)\right|.
\]
\end{stwierdzenie}
\begin{proof}
Niech $A_i$ będzie macierzą przecięcia $i$-tego ramienia.
Jest to macierz trójdiagonalna o przekątnej
$(-a_{i1},\ldots,-a_{ik_i})$ i jedynkach przy
przekątnej. Ustalmy $i$ i oznaczmy długość
ramienia przez $k$. Dla $x\in\mathbb R^k$ mamy
\[
-x^TA_ix=x_1^2+x_k^2+
\sum_{j=1}^{k-1}(x_j-x_{j+1})^2+
\sum_{j=1}^k(a_{ij}-2)x_j^2.
\]
Równość wynika z rozwinięcia kwadratów i zachodzi
również dla $k=1$ (pierwsze dwa składniki dają
$2x_1^2$). Wszystkie składniki są nieujemne,
a ich jednoczesne znikanie wymusza $x=0$.
Zatem $A_i$ jest ujemnie określona.

Oznaczmy przez $d_j$ wyznacznik ogona dodatniej
macierzy $-A_i$ zaczynającego się w pozycji $j$,
przy $d_{k+1}=1$, $d_{k+2}=0$. Z rozwinięcia
wyznacznika wzdłuż pierwszego wiersza, jak w lemacie
\ref{lem:tg-plumbing-chain-det}, mamy
$d_j=a_{ij}d_{j+1}-d_{j+2}$.
Indukcyjnie $d_j>d_{j+1}>0$ dla $1\le j\le k$:
dla $j=k$ mamy $d_k=a_{ik}\ge2>d_{k+1}=1$,
a w kroku indukcyjnym
$d_j\ge2d_{j+1}-d_{j+2}>d_{j+1}$.
Ponadto $\gcd(d_j,d_{j+1})=\gcd(d_{j+1},d_{j+2})=1$,
a ułamek łańcuchowy, przez podstawienie od końca,
wynosi $d_1/d_2$. Wobec dodatniości i względnej
pierwszości otrzymujemy $d_1=\alpha_i$, $d_2=\beta_i$.
Stąd $\det A_i=(-1)^k\alpha_i$, a przez wzór
na element macierzy odwrotnej
\[
(A_i^{-1})_{11}
=\frac{(-1)^{k-1}\beta_i}{(-1)^k\alpha_i}
=-\frac{\beta_i}{\alpha_i}.
\]
Dla $k=1$ licznik jest wyznacznikiem pustej macierzy.

Po usunięciu centralnego wierzchołka pozostały blok
$A=A_1\oplus\cdots\oplus A_r$ jest odwracalny.
Wektor sprzęgający centrum z ramionami ma jedynkę
przy pierwszym wierzchołku każdego ramienia.
Dopełnienie Schura wynosi zatem
$S=b-\sum_i(A_i^{-1})_{11}=b+\sum_i\beta_i/\alpha_i$.
Ponieważ $\det Q_\Gamma=(\prod_i\det A_i)S$,
po wzięciu modułu otrzymujemy tezę.
\end{proof}

\begin{stwierdzenie}[Określoność macierzy gwiaździstej]
\label{st:tg-star-definiteness}
Niech $\Gamma$ będzie ważonym drzewem gwiaździstym
z centralną wagą $b\in\mathbb Z$
i $r\ge1$ ramionami o wagach
$-a_{i1},\ldots,-a_{ik_i}$, gdzie
$a_{ij}\ge2$. Niech
$\alpha_i/\beta_i=[a_{i1},\ldots,a_{ik_i}]^-$,
a $N=1+\sum_{i=1}^rk_i$ oznacza
liczbę wierzchołków.
Połóżmy
$s=b+\sum_i\beta_i/\alpha_i$.
Wtedy:
\begin{enumerate}[label=(\roman*)]
\item jeśli $s<0$, forma $Q_\Gamma$
jest ujemnie określona;
\item jeśli $s=0$, forma jest ujemnie półokreślona,
a jej jądro ma wymiar jeden;
\item jeśli $s>0$, forma ma dokładnie
jedną dodatnią i $N-1$ ujemnych
wartości własnych.
\end{enumerate}
W przyjętej konwencji danych Seiferta
$s=-e_{\mathbb Q}(M)$.
\end{stwierdzenie}

\begin{proof}
Po umieszczeniu centralnego wierzchołka
na końcu macierz ma postać blokową
\[
Q_\Gamma=\begin{pmatrix}
A&v\\v^T&b
\end{pmatrix},\qquad
A=A_1\oplus\cdots\oplus A_r.
\]
Każdy blok $A_i$ jest ujemnie określony
na mocy rachunku sumy kwadratów w dowodzie
stwierdzenia \ref{prop:tg-star-determinant}.
W szczególności $A$ jest odwracalna.

Wykonując odwracalną zmianę bazy
\[
P=\begin{pmatrix}
I&-A^{-1}v\\0&1
\end{pmatrix}
\]
nad $\mathbb R$ otrzymujemy kongruencję
\[
P^TQ_\Gamma P
=\begin{pmatrix}
A&0\\0&b-v^TA^{-1}v
\end{pmatrix}.
\]
Z obliczenia macierzy odwrotnych ramion
w poprzednim dowodzie mamy
$b-v^TA^{-1}v=s$.
Prawo bezwładności Sylvestera daje zatem
$N-1$ ujemnych kierunków z bloku $A$;
ostatni jest odpowiednio ujemny,
zerowy albo dodatni zależnie od znaku $s$.
Wynik wynika bez odwołania do geometrycznej
klasyfikacji rozwarstwień Seiferta.
\end{proof}

\begin{wniosek}[Jądro macierzy przy zerowej liczbie Eulera]
\label{cor:tg-star-nullspace-boundary}
Niech $\Gamma$ będzie skończonym drzewem gwiaździstym
o centralnej wadze $b$ i ramionach o wagach
$-a_{i1},\ldots,-a_{ik_i}$, gdzie $a_{ij}\ge2$.
Niech $Q_\Gamma$ będzie jego całkowitą macierzą przecięcia,
$N=1+\sum_{i=1}^r k_i$ liczbą wierzchołków,
a $\alpha_i/\beta_i=[a_{i1},\ldots,a_{ik_i}]^-$.
Jeżeli $b+\sum_i\beta_i/\alpha_i=0$, to
$\ker(Q_\Gamma:\mathbb Z^N\to\mathbb Z^N)\cong\mathbb Z$
oraz $H_2(\partial P_\Gamma;\mathbb Z)\cong\mathbb Z$.

Dokładniej, jeśli bloki odpowiadające ramionom złożymy
w macierz $A=A_1\oplus\cdots\oplus A_r$,
a $v$ jest wektorem incydencji ramion z centrum,
to przestrzeń zerowa $Q_\Gamma$ nad $\mathbb Q$
jest rozpięta przez
$\bigl(-A^{-1}v,1\bigr)$ w kolejności
(wierzchołki ramion, centrum).
\end{wniosek}
\begin{proof}
Macierz ma postać
$Q_\Gamma=\begin{psmallmatrix}A&v\\v^T&b\end{psmallmatrix}$.
Niech $(y,c)$ należy do jej jądra nad $\mathbb Q$.
Pierwszy wiersz blokowy daje $Ay+vc=0$,
czyli $y=-A^{-1}vc$, ponieważ wszystkie
macierze ramion są nieosobliwe.
Drugi wiersz jest wtedy równoważny
$\bigl(b-v^TA^{-1}v\bigr)c=0$.
Na mocy wzoru na dopełnienie Schura
współczynnik ten wynosi
$b+\sum_i\beta_i/\alpha_i=0$.
Zatem jądro nad $\mathbb Q$ ma wymiar dokładnie jeden
i jest rozpięte przez wskazany wektor wymierny.
Po pomnożeniu tego wektora przez wspólny mianownik
i podzieleniu wszystkich współrzędnych przez ich
największy wspólny dzielnik otrzymujemy
pierwotny wektor całkowity. Niech $w\in\mathbb Z^N$ będzie tak otrzymanym
wektorem pierwotnym. Jeżeli $z\in\mathbb Z^N$
należy do jądra, to $z=t w$ dla $t\in\mathbb Q$.
Z tożsamości Bézouta istnieją $u_j\in\mathbb Z$
takie, że $\sum_j u_jw_j=1$, ponieważ współrzędne
$w$ mają największy wspólny dzielnik jeden.
Wobec tego $t=\sum_j u_jz_j\in\mathbb Z$.
Odwrotnie każda całkowita wielokrotność $w$
należy do jądra, zatem jest ono równe $\mathbb Zw$.
Wreszcie twierdzenie~\ref{thm:tg-plumbing-boundary}
identyfikuje $H_2(\partial P_\Gamma;\mathbb Z)$
z jądrem całkowitej macierzy przecięcia.
\end{proof}

\begin{przyklad}[Relacja między jądrem plumbingowania a torsją]
\label{ex:tg-star-zero-e-kernel}
Dla drzewa z wagą centralną $-1$ i trzema
jednowierzchołkowymi ramionami o wagach $-3$
macierz przecięcia (centrum na pierwszym miejscu)
wynosi
\[
Q=\begin{pmatrix}
-1&1&1&1\\
1&-3&0&0\\
1&0&-3&0\\
1&0&0&-3
\end{pmatrix}.
\]
Mamy $Q(3,1,1,1)^T=0$, a ponieważ
$(3,1,1,1)$ jest wektorem pierwotnym,
otrzymujemy $\ker Q=\mathbb Z(3,1,1,1)^T$.
Na mocy poprzedniego wniosku
$H_2(\partial P_\Gamma;\mathbb Z)\cong\mathbb Z$.
Jednakże pierwsza homologia brzegu
nie jest wolna: z
Przykładu~\ref{ex:tg-seifert-zero-e-torsion}
oraz identyfikacji brzegu
w Twierdzeniu~\ref{thm:tg-star-seifert}
wynika
\[
H_1(\partial P_\Gamma;\mathbb Z)
\cong\mathbb Z\oplus\mathbb Z/3\mathbb Z.
\]
Przykład podkreśla, że kongruencja
macierzy przez \emph{wymierną} zmianę bazy
(np.\ w dowodzie kryterium określoności)
pozwala odczytać jej bezwładność,
ale nie zachowuje na ogół całkowitej grupy
$\operatorname{coker}Q$; do badania torsji
nie wystarcza dopełnienie Schura nad $\mathbb Q$.
\end{przyklad}

\begin{twierdzenie}[Brzeg plumbingowania gwiaździstego]
\label{thm:tg-star-seifert}
Niech $\Gamma$ będzie ważonym drzewem
gwiaździstym z centralną wagą $b\in\mathbb Z$ i ramionami
$(-a_{i1},\ldots,-a_{ik_i})$, gdzie $a_{ij}\ge2$.
Niech $(\alpha_i,\beta_i)$ będą określone powyższymi
ułamkami łańcuchowymi. Wówczas, po ustaleniu opisanych
wyżej orientacji i baz meridian--włókno, brzeg
$\partial P_\Gamma$ jest rozmaitością Seiferta
$M(0;b;(\alpha_i,\beta_i)_{i=1}^r)$.
\end{twierdzenie}
\begin{proof}[Geometryczna redukcja do wypełnień]
\emph{Krok pierwszy: część centralna.}
Brzeg wiązki dyskowej $E(b)\to S^2$ jest okręgową
wiązką nad $S^2$. Usuwamy z bazy $r$ rozłącznych małych dysków,
w których przyłączane są ramiona, oraz jeszcze
jeden dysk niezwiązany z żadnym ramieniem.
Otrzymujemy powierzchnię $B_{0,r+1}$, również
w przypadku $r=0$. Zorientowana wiązka okręgowa
nad tą powierzchnią jest trywialna: powierzchnia
z niepustym brzegiem ma typ homotopijny grafu,
więc jej druga grupa kohomologii całkowitej jest
zerowa, a klasa Eulera klasyfikuje takie wiązki.
Trywializację wybieramy zgodnie z lokalnymi
trywializacjami na dyskach ramion, przenosząc
całą liczbę Eulera na dodatkowy dysk.
Istotnie zmiana trywializacji przez odwzorowanie
$B_{0,r+1}\to S^1$ pozwala zmieniać stopnie
funkcji przejścia na brzegach o dowolne liczby
całkowite o sumie zero: w abelianizacji jedyną
relacją klas brzegowych jest ich suma, więc
tak zadane stopnie określają klasę w $H^1$,
reprezentowaną przez odwzorowanie do $S^1$.
Sumę stopni wyznacza liczba Eulera $b$.
W konwencji nachyleń ustalonej w definicji
danych Seiferta meridian dodatkowego zwykłego
torusa pełnego ma więc klasę
$[\gamma_0]+b[h]$.

\emph{Krok drugi: redukcja ramienia z zachowaniem
jego torusa brzegowego.}
Pracujemy w lokalnych współrzędnych dodatniego
plumbingu z sekcji \ref{sec:tg-plumbing-definition}.
Przy krawędzi oznaczamy przez $\gamma$ dodatnio
zorientowany brzeg powierzchni pozostałej po
wycięciu dysku, a przez $h$ dodatnio zorientowane
włókno okręgowe. Względem wyciętego dysku
krzywa $\gamma$ ma orientację przeciwną do jego
dodatnio zorientowanego brzegu. Zamiana dysku
bazowego z włóknowym $(z,w)\mapsto(w,z)$
daje więc na wspólnym torusie zależności
\[
[\gamma_{\rm in}]=-[h_{\rm par}],\qquad
[h_{\rm in}]=-[\gamma_{\rm par}],
\]
gdzie $\rm in$ odnosi się do strony ramienia,
a $\rm par$ do strony jego poprzednika.
Macierz tej zmiany bazy ma wyznacznik $-1$,
zgodnie z odwróceniem orientacji brzegów przy
sklejaniu. Na każdym wierzchołku normalizujemy
trywializację jak w kroku pierwszym: całkowitą
liczbę Eulera przenosimy na dodatkowy zwykły dysk.

Rozważmy najpierw ostatni wierzchołek o wadze $-a$.
Po usunięciu dysku przy krawędzi pozostaje
wiązka okręgowa nad dyskiem, czyli torus pełny.
Aby obliczyć meridian w ustalonej trywializacji,
usuwamy chwilowo także dodatkowy zwykły dysk.
W iloczynie pierścienia z okręgiem zachodzi
$[\gamma_0]+[\gamma_{\rm in}]=0$,
a zwykłe wypełnienie narzuca relację
$[\gamma_0]-a[h_{\rm in}]=0$.
Stąd meridian ma klasę
$[\gamma_{\rm in}]+a[h_{\rm in}]$.
Po zamianie współrzędnych plumbingu i zmianie
znaku całej klasy otrzymujemy na torusie poprzednika
$a[\gamma_{\rm par}]+[h_{\rm par}]$.
Jest to para $(a,1)$.

W kroku indukcyjnym niech kolejny wierzchołek
o wadze $-a$ ma otwór wejściowy i wyjściowy.
Dalsza część ramienia, zgodnie z założeniem
indukcyjnym, jest torusem pełnym, którego meridian
na torusie wyjściowym ma postać
$\alpha[\gamma_{\rm out}]+\beta[h]$,
gdzie $\alpha>\beta>0$ są względnie pierwsze.
Po zwykłym wypełnieniu dodatkowego dysku
część tego wierzchołka jest wiązką okręgową
nad pierścieniem, czyli $T^2\times[0,1]$.
Doklejenie dalszej części ramienia daje zatem
torus pełny z dołączonym kołnierzem, nie tylko
przestrzeń o tej samej homologii.
W homologii części po zwykłym wypełnieniu mamy
\[
[\gamma_{\rm in}]+[\gamma_{\rm out}]=-a[h].
\]
Po podstawieniu do meridianu dalszej części
dostajemy klasę na torusie wejściowym
\[
-\alpha[\gamma_{\rm in}]+(\beta-a\alpha)[h].
\]
Jest ona pierwotna, gdyż
$\gcd(\alpha,a\alpha-\beta)=\gcd(\alpha,\beta)=1$,
i ogranicza dysk w otrzymanym torusie pełnym:
jest przeniesieniem meridianu przez kołnierz.
Na torusie poprzednika odpowiada jej klasa
\[
(a\alpha-\beta)[\gamma_{\rm par}]
+\alpha[h_{\rm par}].
\]
Otrzymujemy zatem rekurencję par
$(\alpha,\beta)\mapsto(a\alpha-\beta,\alpha)$,
czyli rekurencję ułamka $a-1/(\alpha/\beta)$.
Ponieważ $a\ge2$ oraz $\alpha>\beta>0$,
nowy licznik jest większy od nowego mianownika.

Indukcja od liścia aż do centrum pokazuje,
że całe $i$-te ramię jest torusem pełnym
o meridianie $\alpha_i[\gamma_i]+\beta_i[h]$.
Jest to geometryczna redukcja względna, zachowująca
dane na torusie przyłączenia. Algebraicznie
odpowiada regule \emph{slam-dunk}
z lematu \ref{lem:tg-kirby-slam-dunk},
lecz nie wymaga utożsamiania samej homologii
z typem dyfeomorfizmu ani stosowania lematu
dla izolowanego splotu Hopfa do całego splotu.
Po zredukowaniu wszystkich ramion otrzymujemy
dokładnie konstrukcję z definicji danych Seiferta.
Dla $r=0$ pozostaje samo zwykłe wypełnienie.
\end{proof}

\begin{center}
\begin{tikzpicture}[scale=.95,>=Stealth]
  \node[tgRedVertex,minimum size=11mm] (c) at (0,0) {$b$};
  \node[tgBlueVertex] (a1) at (-1.35,.95) {$-a_{11}$};
  \node[tgBlueVertex] (a2) at (-2.8,.95) {$-a_{12}$};
  \node[tgTealVertex] (b1) at (1.35,.95) {$-a_{21}$};
  \node[tgPurple!80!black,draw=tgPurple,fill=tgPurple!12,circle,
       minimum size=9mm] (d1) at (0,-1.35) {$-a_{31}$};
  \node[tgPurple!80!black,draw=tgPurple,fill=tgPurple!12,circle,
       minimum size=9mm] (d2) at (0,-2.55) {$-a_{32}$};
  \draw[tgAccent] (c)--(a1);
  \draw[tgEdge] (a1)--(a2);
  \draw[tgTree] (c)--(b1);
  \draw[tgEdge] (c)--(d1)--(d2);
  \node[text=tgNavy,font=\footnotesize] at (0,-3.20)
  {Każde ramię zastępujemy jedną parą $(\alpha_i,\beta_i)$.};
\end{tikzpicture}
\end{center}

\begin{przyklad}[Brzeg plumbingowania $E_8$: grupa Poincarégo]
\label{ex:tg-seifert-E8}
Dla $E_8$ centralna waga wynosi $b=-2$,
a trzy ramiona mają odpowiednio długości $1,2,4$
z wagami $-2$. Dlatego
$[2]^-=2/1$, $[2,2]^-=3/2$ i
$[2,2,2,2]^-=5/4$. Na mocy
\ref{thm:tg-star-seifert} brzeg jest rozmaitością
\[
M(0;-2;(2,1),(3,2),(5,4)).
\]
Jej grupa fundamentalna ma prezentację
\[
\left\langle c_1,c_2,c_3,h\ \middle|\
h\text{ centralny},\ c_1^2h=1,\ c_2^3h^2=1,
\ c_3^5h^4=1,\ c_1c_2c_3=h^{-2}\right\rangle.
\]
Po przejściu do ilorazu przez centralną podgrupę
$\langle h\rangle$ otrzymujemy grupę
$\Delta(2,3,5)=\langle x,y,z\mid x^2=y^3=z^5=xyz=1\rangle$.
Aby zobaczyć, że iloraz ten jest nietrywialny,
można wziąć w grupie alternującej $A_5$ permutacje
$s=(12)(34)$ i $t=(235)$. Wtedy $s^2=t^3=(st)^5=1$,
a $s,t$ generują $A_5$, jak wykazano w dowodzie
stwierdzenia \ref{st:tg-E8-poincare}. Przypisania
$x\mapsto s$, $y\mapsto t$, $z\mapsto(st)^{-1}$
określają surjekcję $\Delta(2,3,5)\twoheadrightarrow A_5$.
W szczególności $\pi_1(\partial P_{E_8})\ne1$.
Z poprzedniego przykładu mamy jednocześnie
$H_1(\partial P_{E_8};\mathbb Z)=0$.
To dowodzi, że brzeg jest nietrywialną sferą homologiczną.

Standardowa identyfikacja rozwarstwienia o tych danych
ze sferą Brieskorna $\Sigma(2,3,5)$ noszącą nazwę
\emph{sfery homologicznej Poincarégo} jest osobnym
rezultatem geometrii rozwarstwień Seiferta. Przyjmujemy
tę identyfikację z literatury; obliczenia homologii i
dowód nietrywialności grupy fundamentalnej jej nie
potrzebują. W istocie grupa fundamentalna ma rząd $120$
i jest binarną grupą dwudziestościenną.
\end{przyklad}

\begin{uwaga}[Przypadki z co najwyżej dwoma włóknami wyjątkowymi]
Gdy podstawą jest $S^2$ i liczba włókien wyjątkowych
jest nie większa niż dwa, rozmaitość Seiferta jest
uzyskiwana przez sklejenie dwóch torusów pełnych.
Otrzymujemy przestrzeń soczewkową (w przypadkach
granicznych także $S^1\times S^2$). Wniosek wynika
z rozcięcia $S^2$ wzdłuż okręgu tak, by każdy
z dwóch dysków zawierał co najwyżej jeden punkt
stożkowy (dla zera punktów wybieramy dowolny
okrąg, dla jednego umieszczamy go we wnętrzu
jednego z dysków). Nad dyskiem bez punktów
stożkowych wiązka okręgowa jest trywialna.
Nad dyskiem z jednym punktem stożkowym
otrzymujemy lokalny model wyjątkowego włókna
z dołączonym kołnierzem torusowym, ponieważ
pozostały pierścień nie zawiera punktów
stożkowych. Oba przedobrazy są zatem torusami
pełnymi. Po sklejeniu ich brzegów powstaje
przestrzeń soczewkowa, jeżeli meridiany mają
niezerową algebraiczną liczbę przecięć na torusie;
gdy jest ona zerowa, meridiany wyznaczają tę samą
pierwotną klasę z dokładnością do znaku i wynik
jest dyfeomorficzny z $S^1\times S^2$. Należy odróżnić ten opis od klasyfikacji
orientowanych przestrzeni soczewkowych według
parametrów $(p,q)$, której nie wyprowadzamy tutaj.
\end{uwaga}

\section{Orbicharakterystyka Eulera i geometria}
\label{sec:tg-seifert-orbifold}

\begin{definicja}[Orbicharakterystyka]
Dla orientowalnej orbipowierzchni
$B(g;\alpha_1,\ldots,\alpha_r)$ definiujemy
\[
\chi_{\mathrm{orb}}(B)=2-2g-
\sum_{i=1}^r\left(1-\frac1{\alpha_i}\right).
\]
\end{definicja}

\begin{przyklad}
Dla bazy $S^2(2,3,5)$ mamy
$\chi_{\mathrm{orb}}=2-[(1-1/2)+(1-1/3)+(1-1/5)]
=1/30>0$. Dla trzech punktów stożkowych
wzór upraszcza się do
$\chi_{\mathrm{orb}}=-1+1/\alpha_1+1/\alpha_2+1/\alpha_3$.
Dla $S^2(2,3,6)$ otrzymujemy
$-1+1/2+1/3+1/6=0$, a dla $S^2(2,3,7)$ mamy
$-1+1/2+1/3+1/7=-1/42<0$.
Znaki te odpowiadają odpowiednio trójkątnym
geometriom sferycznej, euklidesowej i hiperbolicznej
orbipowierzchni. Nie określają one samodzielnie geometrii
trójrozmaitości: znaczenie ma również wymierna
liczba Eulera danego rozwarstwienia.
\end{przyklad}

\begin{uwaga}[Miejsce w geometryzacji]
Wśród ośmiu geometrii Thurstona rozmaitości Seiferta
realizują, zależnie od danych i przypadków wyjątkowych,
geometrie $S^3$, $S^2\times\mathbb R$, $\mathbb E^3$,
$\mathrm{Nil}$, $\mathbb H^2\times\mathbb R$
i $\widetilde{\mathrm{SL}}_2(\mathbb R)$.
Dokładniej, dla zamkniętej rozmaitości Seiferta
z orientowalną bazą i danymi jak w tym rozdziale
klasyfikację określa para $\chi_{\mathrm{orb}}(B)$
oraz $e=e_{\mathbb Q}(M)$:
\[
\begin{array}{c|ccc}
 & \chi_{\mathrm{orb}}>0 & \chi_{\mathrm{orb}}=0
 & \chi_{\mathrm{orb}}<0\\ \hline
 e=0 & S^2\times\mathbb R & \mathbb E^3
 & \mathbb H^2\times\mathbb R\\
 e\ne0 & S^3 & \mathrm{Nil}
 & \widetilde{\mathrm{SL}}_2(\mathbb R)
\end{array}
\]
Przyjmujemy to twierdzenie bez dowodu z pracy
P.~Scott, \emph{The geometries of 3-manifolds},
Bull. London Math. Soc. 15 (1983), 401--487,
twierdzenie 5.3(ii), s.~477--478.
Nie jest ono wnioskiem z samego rachunku homologii.

Należy odróżnić geometrię całej rozmaitości
od geometrii konkretnej orbipowierzchni bazowej.
Sfery z jednym punktem stożkowym oraz z dwoma
punktami stożkowymi różnych rzędów są przykładami
złych orbipowierzchni (nie mają nakrycia będącego
rozmaitością; por. Scott, sekcja 2).
Dodatnia orbicharakterystyka nie zapewnia im
geometrii sferycznej. Mimo to rozmaitość Seiferta
nad taką bazą ma geometrię $S^3$:
dla jednego punktu $b+\beta_1/\alpha_1\ne0$,
bo $0<\beta_1/\alpha_1<1$. Dla dwóch punktów
zerowanie wymagałoby $b=-1$ i
$\beta_1/\alpha_1=(\alpha_2-\beta_2)/\alpha_2$.
Oba ułamki są nieskracalne, więc ich równość
wymusza $\alpha_1=\alpha_2$, wbrew założeniu.
W tych przypadkach zatem $e\ne0$ i stosuje się
odpowiednie pole powyższej klasyfikacji.
\end{uwaga}

